\subsection{Degree zero from the two exact sequences and a free coefficient lattice}
\label{subsec:faithful-character-quotient-degree-zero}

This subsection follows
\srcid{SplitHermitianGeneratorsExist.deg\_quot\_eq\_zero}.
The coefficient algebra need not be reduced or etale;
the final step uses its actual ideal filtration.

\begin{srproposition}\label{prop:faithful-character-quotient-degree-zero}
Let \(E\) be a number field with the specified involution
on \(O=\mathcal O_E\), let \(K\) have characteristic
the prime \(p\), and let \(C\) be a source arithmetic
curve over \(K\).
Let \(D\) be the source's
\(\mathrm{ConjData}(C,p,O,d)\), so its lattice bundle
\(\mathcal H\), its tangent lattice \(\mathcal L\),
its actual action \(\rho_H\), and its two exact
sequences are exactly (FC6)--(FC7).
Let \(\chi:O\to K\) be a ring character.
Retain the source's splitting hypothesis: for every
\(a\in O\), the characteristic polynomial of
multiplication by \(a\) on the integral \(O\)-lattice,
mapped to \(K[T]\), splits over \(K\).

Assume there is a fixed \(r\ge0\) such that for every
valuation index \(i\) there are \(x_{i,1},\ldots,x_{i,r}\)
with
\[
 \bigl\{\rho_H(b_j)x_{i,l}:
       1\le j\le[E:\mathbb Q],\ 1\le l\le r\bigr\}
 \quad\text{a }K\text{-basis of }H_K
\]
whose \(R_i\)-span is exactly \(\mathcal H_i\);
here \(b_j\) is the specified integral basis of \(O\).
Then, for the source quotient lattice,
\[
 \deg_C\!\left(\mathcal H/W_\chi\right)=0,\qquad
 W_\chi=\sum_{a\in O}\operatorname{im}
                                  (\rho_H(a)-\chi(a)I).
 \tag{FZ1}\label{eq:faithful-character-quotient-zero}
\]
\end{srproposition}

\paragraph{Degree conventions and two elementary bundle facts.}
For a subspace \(W\subset H_K\), the sublattice is
\(\mathcal H_i\cap W\) and the quotient lattice is
the image of \(\mathcal H_i\) in \(H_K/W\).
These are the source's \(\mathrm{sub}\) and
\(\mathrm{quot}\); no equality of a coefficient ideal's
integral image with its saturation is assumed.

First, degree is additive in an exact lattice sequence.
At each valuation, choose an integral basis of the
submodule and lift a basis of the quotient.
The determinant matrix is block triangular, so its
valuation is the sum of the two determinant valuations.
The product formula removes changes of generic basis.
This proves additivity with the source's weighted
integer-valued degree.

If compatible idempotent endomorphisms act on an exact
lattice sequence and preserve its middle lattice, their
image sequences are also exact lattices.
For a kernel vector, take an arbitrary preimage in the
first term and apply its idempotent; for a quotient
vector, take an integral lift in the middle and apply
the middle idempotent.
The maps commute, so these are the required preimages
and lifts in the image subspaces.
Consequently their degrees add.
This is the source's
\srcid{FilteredExactSequenceProjectorBundles.subDeg\_range\_shortExact}.

Second, Frobenius multiplies subbundle degree by \(p\):
\[
       \deg_C\bigl(F^*\mathcal V\cap F(W)\bigr)
                               =p\,\deg_C(\mathcal V\cap W).
 \tag{FZ2}\label{eq:faithful-frob-subdegree}
\]
Here \(F(W)\) is the \(K\)-span of coordinatewise
\(p\)-th powers.
The full-lattice formula is
\srcid{FrobeniusLatticeDegreeMulP.deg\_frobLat}; the required
subspace formula is separately
\srcid{LinearIndependentMapVec.SubFrob.subDeg\_frob}.
Here is its intersection argument.
Choose a full integral local basis \(b\) adapted to \(W\),
whose first \(m=\dim W\) elements form an integral
basis \(c\) of the intersection.
The coordinatewise \(p\)-th powers of \(b\) are a full
basis, because its determinant is raised to its
\(p\)-th power.
The \(p\)-th powers of \(c\) form a basis of the
Frobenius span of \(W\).
More generally the source proves their linear independence
by taking a left inverse to their original row matrix
and applying Frobenius to the matrix identity; it does
not take \(p\)-th roots of coefficients.
The local Frobenius lattice is the integral span of
the new full basis: inclusion of all Frobenius
generators follows from their original integral
coordinates, and the reverse inclusion uses the
new basis vectors themselves.
In those coordinates, its intersection with the
new \(m\)-dimensional subspace is precisely the span
of the first \(m\) vectors.
Thus it has the asserted transformed intersection basis.
The semilinear coordinate formula sends every
change-of-basis matrix \(M\) on \(W\) to \(M^{(p)}\);
hence \(\det(M^{(p)})=\det(M)^p\).
Valuations multiply by \(p\), proving (FZ2).
The full-lattice finite-exception condition gives the
corresponding lattice bundles throughout.
No surjectivity of Frobenius on \(K\) is asserted.

\paragraph{Finite constants and their action.}
The kernel of \(\chi\) contains the nonzero element
\(p\in O\).
Every nonzero ideal of \(O\) has finite quotient:
if \(a\ne0\) belongs to it, multiplication by \(a\)
has nonzero integral determinant on the integral basis
of \(O\), so \(O/aO\), and hence the ideal quotient,
is finite.
Thus \(\chi(O)\) is a finite integral domain.
For every nonzero element, multiplication is injective
and therefore surjective on this finite set, so it has
an inverse there.
It is a finite subfield \(k\subset K\).
Write \(|k|=p^s\), with \(s>0\), and
\(\chi_0:O\to k\) for the same character with restricted
codomain.

Every nonzero element of \(k\) has finite multiplicative
order; hence every valuation \(v_i\) of \(K\) is \(1\)
on it.
In particular \(k\subset R_i\) for all \(i\).
Set
\[
 A=k\otimes_{\mathbb Z}O,\quad
 \chi_A(c\otimes a)=c\,\chi_0(a),\quad
 \mathfrak m=\ker\chi_A .
\]
The algebra \(A\) is a finite ring.
Extend the three \(O\)-actions on
\(\mathcal H,\mathcal L,\mathcal L^\vee\) to \(A\)
by \(c\otimes a\mapsto c\,\rho(a)\).
They preserve their lattices because \(c\in R_i\).
The same is true of every idempotent in \(A\).

\paragraph{Frobenius forces every constant idempotent degree to vanish.}
For \(e^2=e\in A\), let
\[
                            d_H(e)=\deg_C(\mathcal H\cap eH_K).
\]
The idempotents induced on the first and last Hodge
terms give, by exactness,
\[
 d_H(e)=
 \deg_C\bigl(\mathcal L^\vee\cap e_{\rm dual}L_K^\vee\bigr)
 +\deg_C\bigl(\mathcal L\cap e_{\rm tan}L_K\bigr).
 \tag{FZ3}\label{eq:faithful-idempotent-Hodge-degree}
\]
The action on the dual is the specified conjugate
transpose action, not an unconjugated replacement.

Define the source's coefficient Frobenius
\[
                \mathfrak F:A\to A,\qquad
                             c\otimes a\longmapsto c^p\otimes a.
\]
The action of \(\mathfrak F(e)\) on either conjugate
term of (FC6) is the coordinatewise Frobenius of the
action of \(e\) on the corresponding Hodge term.
Explicitly, for any matrix representation \(\rho\),
\[
 \operatorname{act}_{\rho^{(p)}}(\mathfrak F(e))
                              =\operatorname{act}_\rho(e)^{(p)}.
\]
The range of the matrix on the right is the Frobenius
span of its old range: one inclusion follows by
taking \(p\)-th powers, and the reverse follows by
applying the matrix to standard basis vectors,
which are their own \(p\)-th powers.
Applying the exact conjugate sequence, then (FZ2)
to its two terms, and comparing with (FZ3), yields
\[
                              d_H(\mathfrak F(e))=p\,d_H(e).
 \tag{FZ4}\label{eq:faithful-idempotent-Frobenius-step}
\]
This is the source's \srcid{step}.

Coefficient Frobenius has period \(s\) on \(A\), because
\(c^{p^s}=c\) on \(k\).
Its iterates preserve idempotents.
Iterating (FZ4) therefore gives
\[
                          d_H(e)=p^s d_H(e).
\]
The degree is an integer and \(p^s>1\), so
\[
                                  d_H(e)=0.
 \tag{FZ5}\label{eq:faithful-idempotent-degree-zero}
\]
This is \srcid{vanish}.
It has used the two exact sequences and finite constants,
not the free coefficient-lattice premise yet.

\paragraph{The source's idempotent for the character block.}
For every element \(t\) of a finite ring there is a
positive integer \(N_t\) with \(t^{N_t}\) idempotent.
Indeed a repeated pair of powers gives an eventually
periodic sequence; choose a sufficiently large positive
multiple of that period.
Define exactly as in \srcid{exists\_idem}
\[
                          e=\prod_{t\in\mathfrak m}(1-t^{N_t}).
 \tag{FZ6}\label{eq:faithful-character-idempotent}
\]
Each factor is idempotent, and the factors commute;
hence \(e^2=e\).
Modulo \(\mathfrak m\), every factor is \(1\), so
\(\chi_A(e)=1\).
For \(t\in\mathfrak m\), the factor \(1-t^{N_t}\)
shows that \(e t\) is nilpotent.
In fact
\(t^{N_t}(1-t^{N_t})=0\), and taking the corresponding
power of the product kills \((e t)^{N_t}\).

These statements do not require \(A\) to be semisimple.
They supply the actual idempotent and nilpotence
hypotheses used in the source's next step.

\paragraph{A one-dimensional descending ideal quotient.}
For any nonzero ideal \(J\subset eA\), the source
constructs an ideal \(J_1<J\), an element \(t\in J\),
and a \(k\)-linear functional \(l:A\to k\), such that
\[
 \begin{gathered}
       J=J_1+kt,\qquad \mathfrak m t\subset J_1,\\
                         l|_{J_1}=0,\qquad l(t)=1.
 \end{gathered}
 \tag{FZ7}\label{eq:faithful-ideal-colength-one}
\]
Here is its determinant-trick argument, with the
nilpotence premise retained.
Since \(A\) is finite, \(J\) is a finitely generated
\(A\)-module.
If \(J\subset\mathfrak mJ\), the determinant trick gives
an element \(r\equiv1\pmod{\mathfrak m}\) annihilating \(J\).
Then \(1-r\in\mathfrak m\) acts as the identity on \(J\).
Also \(e\) acts as the identity on \(J\subset eA\).
Thus \(e(1-r)\) acts as the identity on \(J\), although
it is nilpotent by the preceding paragraph.
This forces \(J=0\), a contradiction.
Hence choose \(t\in J\setminus\mathfrak mJ\).

Choose a \(k\)-linear functional on \(A/\mathfrak mJ\)
nonzero on the class of \(t\), and scale it to value \(1\).
Its pullback is \(l\).
Put \(J_1=J\cap\ker l\).
For \(a\in A,b\in J\), write
\[
                  ab=\chi_A(a)b+(a-\chi_A(a))b.
\]
The second term belongs to \(\mathfrak mJ\), on which
\(l\) vanishes.
Therefore \(l(ab)=\chi_A(a)l(b)\), proving \(J_1\)
is an ideal.
The equations in (FZ7) follow by subtracting \(l(b)t\)
from any \(b\in J\), and by
\(\mathfrak m t\subset\mathfrak mJ\).
This is precisely
\srcid{IdealDescendingOneDimQuotient.exists\_colength\_one}.

\paragraph{The local free coordinates in the filtration step.}
Write \(\tau:A\to\operatorname{End}_K(H_K)\) for the
extended action, and set
\[
       W_J=\sum_{a\in J}\operatorname{im}\tau(a),\qquad
       Q=\mathcal H/W_{\mathfrak m}.
\]
The additional hypothesis in the proposition now enters.
The given local vectors \(x_{i,l}\), together with the
integral basis \(b_j\), give a free basis over
\(R_i\otimes_k A\) of \(\mathcal H_i\).
For a \(k\)-linear functional \(l:A\to k\), the resulting
coefficient functionals
\(\Theta_{i,l,\nu}:H_K\to K\) satisfy
\[
 \Theta_{i,l,\nu}(\tau(a)x_{i,\mu})
                               =\delta_{\nu\mu}\,l(a).
\]
They are integral on \(\mathcal H_i\):
in the specified integral basis the coefficients
\(l(1\otimes b_j)\) lie in \(k\subset R_i\).
This is the source's \(\mathrm{FreeAt}\) construction;
the local basis is allowed to depend on \(i\).

For the \(l,t,J_1,J\) in (FZ7), define
\[
                           P_i(y)=
                           \sum_\nu\Theta_{i,l,\nu}(y)x_{i,\nu}.
\]
It is \(K\)-linear and preserves \(\mathcal H_i\).
Expansion in the local free coefficient basis gives
three precise properties:
\[
 \begin{split}
 y-\tau(t)P_i(y)&\in W_{J_1} &&(y\in W_J),\\
                         P_i(W_{J_1})&=0,\\
 P_i(\tau(t)z)-z&\in W_{\mathfrak m} &&(z\in H_K).
 \end{split}
 \tag{FZ8}\label{eq:faithful-filtration-local-maps}
\]
For the first line, apply
\(a-l(a)t\in J_1\) to every coefficient \(a\in J\).
For the second, \(l\) vanishes on the ideal \(J_1\).
For the third,
\(ta-\chi_A(a)t\in J_1\), so
\(l(ta)=\chi_A(a)\); subtracting \(\chi_A(a)\)
from each coefficient leaves an element of
\(\mathfrak m\).
These are the coordinate calculations in
\srcid{AlgebraActionBundleFiltrationStep.exists\_local}.

If the valuation index set is empty, every degree is
already zero, and (FZ1) follows.
Otherwise fix one index \(i_0\).
Define a global \(K\)-linear map
\[
                         g:W_J\longrightarrow H_K/W_{\mathfrak m},
                   \qquad g(y)=[P_{i_0}y].
\]
Equations (FZ8) prove that its kernel is \(W_{J_1}\)
and that \(g(\tau(t)z)=[z]\).
They also give exactness at every lattice.
For an integral class \([z]\), use the integral lift
\(\tau(t)z\in\mathcal H_i\cap W_J\).
Conversely, if \(y\in\mathcal H_i\cap W_J\), then
\[
                              [P_{i_0}y]=[P_i y].
\]
To verify this equality, use
\(y-\tau(t)P_i y\in W_{J_1}\), apply \(P_{i_0}\),
and then use the last line of (FZ8).
The right side is an integral quotient class because
\(P_i\) preserves \(\mathcal H_i\).
Thus the source's exact lattice sequence is
\[
 0\longrightarrow \mathcal H\cap W_{J_1}
  \longrightarrow \mathcal H\cap W_J
  \stackrel g\longrightarrow Q\longrightarrow0.
\]
Degree additivity gives
\[
      \deg_C(\mathcal H\cap W_J)
          =\deg_C(\mathcal H\cap W_{J_1})+\deg_C Q.
 \tag{FZ9}\label{eq:faithful-ideal-degree-step}
\]
This proof has used the actual local coordinate maps;
it has not assumed that the same coefficient basis
works at every valuation.

\paragraph{Completion of the proof.}
The ideal set of the finite ring \(A\) is finite.
Apply (FZ7)--(FZ9) by descending induction, starting
with \(J=eA\) and ending at \(0\).
As \(\chi_A(e)=1\), the starting ideal is nonzero.
The source therefore obtains a positive integer \(N\)
with
\[
                              d_H(e)=N\deg_C Q.
\]
By (FZ5) its left side is zero; hence \(\deg_C Q=0\).

Finally,
\[
 W_{\mathfrak m}
    =\sum_{a\in O}\operatorname{im}
                              (\rho_H(a)-\chi(a)I).
\]
One inclusion follows because
\(1\otimes a-\chi_0(a)\otimes1\in\mathfrak m\).
For the reverse inclusion, modulo these elements every
\(c\otimes a\) equals \(c\chi_0(a)\otimes1\);
an element of \(\mathfrak m\) then becomes zero.
Thus the ideal is generated by those differences,
and its action image is the displayed sum.
The quotient \(Q\) is exactly the quotient in (FZ1),
proving the proposition.

The source statement retains the splitting hypothesis
on the multiplication characteristic polynomials.
The displayed source proof body does not use that
hypothesis; it is nevertheless retained above, so the
statement has not been generalized here.

\paragraph{Frozen-source ledger and actual geometric use.}
The final source is
\srcid{src/Development/AutomorphicSpectrum/Proofs/AlgebraProofs/SplitHermitianGeneratorsExist.lean}
(blob \srcid{659a82dce2a636fc33707996c8856ebc090a819b}).
The actual supporting producers are, in that same directory:
\begin{itemize}
\item \srcid{FilteredExactSequenceProjectorBundles.lean},
blob \srcid{5a1e6aefb1675481a726fef2ff75faddb1820c28};
\item \srcid{FrobeniusLatticeDegreeMulP.lean},
blob \srcid{e56760012070cfb0458badd029e0b0742304f85e};
\item \srcid{LinearIndependentMapVec.lean},
blob \srcid{2f83f8d798ab9f9e4371667c6ffbf32d4902f908},
especially \srcid{SubFrob.subDeg\_frob}; this is the
actual subspace/intersection Frobenius producer;
\item \srcid{ModuleBasisLatticeBundleFree.lean},
blob \srcid{9753258f9343bf632b4b9784f4e89d3e18468af4};
\item \srcid{IdealDescendingOneDimQuotient.lean},
blob \srcid{3980cc27b82f3dbd9e6ba1b42081c6cc96dd2f2d};
\item \srcid{AlgebraActionBundleFiltrationStep.lean},
blob \srcid{757a0401673a313b13c7c87787424ea88c678a00}.
\end{itemize}
All were checked against the frozen baseline tag.

For an actual vertical curve in a supplied model,
Proposition~\ref{prop:faithful-actual-family-ex} and
Proposition~\ref{prop:faithful-family-two-exact}
supply these exact sequences with their actual tangent
and full formal lattices, once the model has
\(\mathrm{PowFns}\) and the specified \(\mathrm{FibPF}\)
at the curve centres.
If the actual formal lattices also have the fixed-rank
free \(O\)-bases stipulated in this proposition,
(FZ1) proves degree zero for their actual character
quotient determinant.
This is the source's degree-zero conclusion on the
vertical curve.  The free-basis input is precisely
where the separate formal-kernel/freeness producer
must be joined; it has not been replaced by ordinary
\(R_i\)-freeness or by the existence of the two sequences.

This conclusion does not construct the full Mumford
model, identify its analytic flat frame with a canonical
determinant extension, or pull arbitrary eigensaturated
lattices through a nonflat base change.
Those are separate joins in the flat-saturation producer.
